%%%%%%%%%%%%%%%%%%%%%%%%%%%%%%%%%%%%%%%%%%%%%%%%%%%%%%%%%%%%%%%%%%%%%%%%%%%%%%%
% 
% Semi-Analytic Sampling from an SCF Potential
% For the DiscO project
% 
%%%%%%%%%%%%%%%%%%%%%%%%%%%%%%%%%%%%%%%%%%%%%%%%%%%%%%%%%%%%%%%%%%%%%%%%%%%%%%%

%%%%%%%%%%%%%%%%%%%%%%%%%%%%%%%%%%%%%%%%%%%%%%%%%%%%%%%%%%%%%%%%%%%%%%%%%%%%%%%
%%% SETUP

\documentclass[12pt, english]{article}

%%%%%%%%%%%%%%%%%%%%%%%%%%%%%%%%%%%%%%%%%%%%%%%%%%%%%%%%%%%%%%%%%%%%%%%%%%%%%%%
%                                 PACKAGES
%%%%%%%%%%%%%%%%%%%%%%%%%%%%%%%%%%%%%%%%%%%%%%%%%%%%%%%%%%%%%%%%%%%%%%%%%%%%%%%
% Importing Packages

% Page Layout
\usepackage{fancyhdr}
\usepackage{geometry}

% Figures
\usepackage{graphicx}

% Citations
\usepackage{natbib}

% Math
\usepackage{amsmath,amstext,amssymb}
\usepackage{dsfont}


%%%%%%%%%%%%%%%%%%%%%%%%%%%%%%%%%%%%%%%%%%%%%%%%%%%%%%%%%%%%%%%%%%%%%%%%%%%%%%%
%                                 SETTINGS
%%%%%%%%%%%%%%%%%%%%%%%%%%%%%%%%%%%%%%%%%%%%%%%%%%%%%%%%%%%%%%%%%%%%%%%%%%%%%%%
% Implementing Package Settings

\graphicspath{{images/}{../images/}}

\geometry{total={210mm,297mm},
          left=25mm,right=25mm,
          top=20mm,bottom=20mm,
          bindingoffset=0mm}

\raggedbottom
\raggedright


%%%%%%%%%%%%%%%%%%%%%%%%%%%%%%%%%%%%%%%%%%%%%%%%%%%%%%%%%%%%%%%%%%%%%%%%%%%%%%%
%                                  DOCUMENT
%%%%%%%%%%%%%%%%%%%%%%%%%%%%%%%%%%%%%%%%%%%%%%%%%%%%%%%%%%%%%%%%%%%%%%%%%%%%%%%

\begin{document}


% -----------------------------------------------------------------------------
%                                  TITLE PAGE
% -----------------------------------------------------------------------------

\pagestyle{fancy}
\lhead{Semi-Analytic Sampling from an SCF Density}
\rhead{DiscO $|$ \thepage}
\cfoot{Contributions to Date: Nathaniel Starkman, Jo Bovy, Chris Carr} % get rid of the bottom page number

% -----------------------------------------------------------------------------
%                                     TOC
% -----------------------------------------------------------------------------

% \tableofcontents
% \let\tableofcontents\relax
% \listoffigures
% \listoftables

% -----------------------------------------------------------------------------
%                                  CONTENT
% -----------------------------------------------------------------------------

The current Galpy implementation of the SCF density field. This is very similar to other implementations, different mostly in whether the full $Y_{lm}$ is used over the associated Legendre.

\begin{flalign}
    \rho(r, \theta, \phi) &= \sum_{n,l=0}^{\infty} \sum_{m=0}^{l} \tilde{\rho}_{nlm}(r) \left[A_{nlm} \cos({m\phi}) + B_{nlm} \sin{m\phi} \right] P_{l}^{m}(\cos{\theta}) \label{eq:rho} &
\end{flalign}

\textbf{Note I made a mistake in $\tilde{\rho}_{nlm}(r)$. I think it is actually in 2 indices $\tilde{\rho}_{nl}(r)$. This doesn't make a difference, except to clean up some math here or there.}

\section{Sampling from CDFs} % (fold)
\label{sec:sampling_from_cdfs}


    \subsection{CDF(r)} % (fold)
    \label{sub:cdf_r}

        The CDF in $r$ is

        \begin{flalign}
            \mathbf{R}(r) &= \frac{1}{N} \int_0^r dV \rho(r', \theta', \phi') = \frac{m(\leq r)}{M_{tot}} &
        \end{flalign}

        The numerator and denominator are galpy functions.

        Numerically sampling from the CDF is challenging in these coordinates because the domain ($r$) is infinite on the half-open interval $[0, \infty)$. We switch to the a finite domain

        \begin{equation}\label{eq:ho_2_15}
         \zeta = \frac{r-1}{r+1}
        \end{equation}

        Where the points $r=0, \infty$ correspond to $\zeta=-1, 1$, respectively.

        After a set of sample points $\{\zeta\}$ is drawn, they are transformed back to $r$ using

        \begin{equation}\label{eq:ho_2_16}
         r = \frac{1+\zeta}{1-\zeta}
        \end{equation}


    % subsection cdf_r (end)

    \subsection[CDF(theta;r)]{CDF($\theta$;r)} % (fold)
    \label{sub:cdf_theta}

        First lets start with a definition: $\theta \in [-\frac{\pi}{2}, \frac{\pi}{2}]$, from the South to North pole, respectively.

        \textbf{Double check this is compatible with galpy's definition. Will require mapping to array indices, otherwise.}
        \vspace{10pt}

        The CDF for $x=\cos{\theta}$, at a given $r$ is

        \begin{flalign}
            \mathbf{\Theta}(\cos^{-1}x; r)
            &= \frac{1}{N} \int_{0}^{2\pi}\rm{d}\phi' \int_{0}^{\cos^{-1}x}\rm{d}\theta' r^2 \sin(\theta') \rho(r', \theta', \phi') & \\
        \intertext{\quad We first change to $\rm{d}\cos(\theta) = -\sin(\theta)\rm{d}\theta$.}
            &= \frac{r^2}{N} \int_{-1}^{x}\rm{d}\cos\theta' \int_{0}^{2\pi}\rm{d}\phi' \; \rho(r', \theta', \phi') & \\
        \intertext{\quad Substituting in \eqref{eq:rho}. Note that over the interval $[0, 2\pi)$ only $m=0$ for $\cos{m\phi}$ is nonzero. Therefore can can also simplify to find}
            &= \frac{r^2}{N} \int_{-1}^{x}\rm{d}\cos\theta' \int_{0}^{2\pi}\rm{d}\phi' \; \sum_{n,l=0}^{n,l_{\max}} \tilde{\rho}_{nl0}(r) A_{nl0} P_{l}^{0}(\cos\theta') & \\
        \intertext{\quad No term depends on $\phi$, so we do the integral}
            &= \frac{2\pi r^2}{N} \int_{-1}^{x}\rm{d}\cos\theta' \sum_{n,l=0}^{n,l_{\max}} \tilde{\rho}_{nl0}(r) A_{nl0} P_{l}^{0}(\cos\theta') & \\    
        \intertext{\quad Simplifying, moving in the integral}
            &= \frac{2\pi r^2}{N} \sum_{n,l=0}^{n,l_{\max}} \tilde{\rho}_{nl0}(r) A_{nl0} \int_{-1}^{x}\rm{d}\cos\theta' \; P_{l}^{0}(\cos{\theta'}) & \\
        \intertext{\quad We can do the $n$-sum. Let $Q_l(r) = \sum_{n=0}^{n_{\max}}A_{nl0} \tilde{\rho}_{nl0}(r)$ }
            &= \frac{2\pi r^2}{N} \sum_{l=0}^{l_{\max}} Q_l(r) \int_{-1}^{x}\rm{d}\cos\theta' \; P_{l}^{0}(\cos{\theta'}) & \\
        \intertext{\quad Doing the integral}
            &= \frac{2\pi r^2}{N} \sum_{l=0}^{l_{\max}} Q_l(r) \left[\frac{2\sin(\pi l)}{\pi l (l+1)} + \frac{\sqrt{1-x^2} P_{l+1}^{1}(x)}{l(l+1)} \right] & \\
        \intertext{Note that since $l\in \mathds{Z}$, only $l=0$ is nonzero in the $\sin$ term.}
            &= \frac{2\pi r^2}{N} \sum_{l=0}^{l_{\max}} Q_l(r) \left[2\delta_{l0} + \frac{\sqrt{1-x^2} P_{l+1}^{1}(x)}{l(l+1)} \right] \\
            &= \frac{2\pi r^2}{N} \left[2 Q_0(r) + \sum_{l=0}^{l_{\max}} Q_l(r) \frac{\sqrt{1-x^2} P_{l+1}^{1}(x)}{l(l+1)} \right] \label{eq:theta_unnormalized}
        % \intertext{\quad Doing the integral}
        %     &= \frac{2\pi r^2}{N} \sum_{n,l=0}^{n,l_{\max}} A_{nl0} \tilde{\rho}_{nl0}(r) \left[\frac{2\sin(\pi l)}{\pi l (l+1)} + \frac{P_{l+1}(x) - P_{l-1}(x)}{2l + 1} \right] & \\
        % \intertext{Note that since $l\in \mathds{Z}$, only $l=0$ is nonzero in the $\sin$ term.}
        %     &= \frac{2\pi r^2}{N} \sum_{n,l=0}^{n,l_{\max}} A_{nl0} \tilde{\rho}_{nl0}(r) \left[2\delta_{l0} + \frac{P_{l+1}(x) - P_{l-1}(x)}{2l + 1} \right] & \\
        % \intertext{\quad We can do the $n$-sum. Let $Q_l(r) = \sum_{n=0}^{n_{\max}}A_{nl0} \tilde{\rho}_{nl0}(r)$ }
        %     &= \frac{2\pi r^2 }{N} \sum_{l=0}^{l_{\max}} Q_l(r) \left[2\delta_{l0} + \frac{P_{l+1}(x) - P_{l-1}(x)}{2l + 1} \right]  \label{eq:theta_unnormalized}
        \end{flalign}

        The normalization constant is the above, but with $x \rightarrow 1$. The last term in \eqref{eq:theta_unnormalized} is 0 for $x=1$, leaving
        $N = 4\pi r^2 Q_0(r)$.

        Therefore,

        \begin{flalign}
            \mathbf{\Theta}(x; r)
            &= 1 + \frac{1}{2 Q_0(r)}\sum_{l=0}^{l_{\max}} Q_l(r) \frac{\sqrt{1-x^2} P_{l+1}^{1}(x)}{l(l+1)}  \\
            &= \frac{1+x}{2} + \frac{1}{2 Q_0(r)}\sum_{l=1}^{l_{\max}} Q_l(r) \frac{\sqrt{1-x^2} P_{l+1}^{1}(x)}{l(l+1)}
        \intertext{\quad an equivalent form, using a recurrence relation:}
            &= \frac{1+x}{2} + \frac{1}{2 Q_0(r)} \sum_{l=1}^{l_{\max}} Q_l(r) \frac{P_{l+1}(x) - P_{l-1}(x)}{2l + 1}
        \end{flalign}

        As a computational note, many terms can be pre-computed. Also if a minimum resolution is set, than many terms can be computed on a grid and interpolated / splined. This will prevent computational time from scaling with the number of sample points.

        $=>$ \textbf{lay grid in r and theta and spline this CDF for use if sampling many points.}    

    % subsection cdf_theta (end)

    \subsection[CDF(phi;r,theta)]{CDF($\phi$;r,$\theta$)} % (fold)
    \label{sub:cdf_phi}

        Let $y = \phi \in [0, 2\pi)$

        \begin{flalign}
            \mathbf{\Phi}(y; r, \theta)
            &= \frac{1}{N} \int_{0}^{y} \rm{d}\phi \; r\sin\theta \sum_{n,l=0}^{n/l_{\max}} \sum_{m=0}^{l} \tilde{\rho}_{nlm}(r) \left[A_{nlm} \cos({m\phi}) + B_{nlm} \sin{m\phi} \right] P_{l}^{m}(\cos{\theta}) & \\
        \intertext{Pulling in the integral}
            &= \frac{r\sin\theta}{N} \sum_{n,l=0}^{n/l_{\max}} \sum_{m=0}^{l} \tilde{\rho}_{nlm}(r) P_{l}^{m}(\cos{\theta}) \int_{0}^{y} \rm{d}\phi \left[A_{nlm} \cos({m\phi}) + B_{nlm} \sin{m\phi} \right] & \\
        \intertext{Performing the integrals}
            &= \frac{r\sin\theta}{N} \sum_{n,l=0}^{n/l_{\max}} \sum_{m=0}^{l} \tilde{\rho}_{nlm}(r) P_{l}^{m}(\cos{\theta}) \left[A_{nlm} \frac{\sin(m y)}{m} + B_{nlm} \frac{1-\cos(m y)}{m} \right] & \\
        \intertext{The $n$-sum does not depend on $y$, so let $[R/S]_{l}^{m}(r,x)= \left(\sum_{n=0}^{n_{\max}} [A/B]_{nlm} \tilde{\rho}_{nlm}(r) \right) \sqrt{1-x^2} P_{l}^{m}(x)$}
            &= \frac{r}{N} \sum_{l=0}^{l_{\max}} \sum_{m=0}^{l} \left[R_{l}^{m}(r,x) \frac{\sin(m y)}{m} + S_{l}^{m}(r,x) \frac{1-\cos(m y)}{m} \right] & \\
        \intertext{
            we are interested in doing the $l$-sum, but cannot since $m$ depends on $l$. By reindexing, this can be done. Let $m'\in [0, l_{\max}] \in \mathds{Z}$ and $l' \in [m', l_{\max}] \in \mathds{Z}$. Now the summation order can be reversed.
        }
            &= \frac{r}{N} \sum_{m'=0}^{l_{\max}} \sum_{l'=m'}^{l_{\max}} \left[R_{l'}^{m'}(r,x) \frac{\sin(m' y)}{m'} + S_{l'}^{m'}(r,x) \frac{1-\cos(m' y)}{m'} \right] & \\
        \intertext{
            Let $[R/S]^{m}(r, x) = \sum_{l=m}^{l_{\max}} [R/S]_{l}^{m}(r,x)$
        }
            &= \frac{r}{N} \sum_{m'=0}^{l_{\max}} \left[R^{m'}(r,x) \frac{\sin(m' y)}{m'} + S^{m'}(r,x) \frac{1-\cos(m' y)}{m'} \right] \\
        \intertext{Note that for $m'=0$ many of these terms disappear}
            &= \frac{r}{N} \left(R^{0}(r,x) \, y + \sum_{m'=1}^{l_{\max}} \left[R^{m'}(r,x) \frac{\sin(m' y)}{m'} + S^{m'}(r,x) \frac{1-\cos(m' y)}{m'} \right] \right) \\
        \end{flalign}

        The normalization constant is the above, but with $y \rightarrow 2\pi$. For $m> 0$ all terms are 0, leaving

        \begin{flalign}
            N &= 2\pi r R^{0}(r,x)
        \end{flalign}


        Therefore

        \begin{equation}
            \mathbf{\Phi}(y; r, \theta) = \frac{y}{2\pi} + \frac{1}{R^{0}(r,x)} \sum_{m'=1}^{l_{\max}} \left[R^{m'}(r,x) \frac{\sin(m'y)}{2\pi m'} + S^{m'}(r,x) \frac{1-\cos(m' y)}{2\pi m'} \right]
        \end{equation}

        Computational notes: the R/S terms can be precomputed along the $m$ axis using a cumulative sum (the vector of summations will be length $m$-to-$l_{\max}$), then multiplied by the $m$-to-$l_{\max}$ vector of $\sin,\cos (y)$.

    % subsection cdf_phi (end)

% section sampling_from_cdfs (end)


\section[Choosing Nmax, Lmax]{Choosing $N_{\max}$, $L_{\max}$} % (fold)
\label{sec:choosing_nmax_lmax}

    An important consideration for basis-function expansion (BFE) representations is that of resolution -- to able to localize a features to a region of volume $V$. BFE representations with low numbers of coefficients
    tend to have quite poor resolution and increasing the number of coefficients will similarly increase resolution. However computations can become quite resource intensive on large expansions.
    The question then, is what is the minimum number of coefficients necessary to achieve a specific resolution?

    In the context of the Self-Consistent Field, this question is posed instead as the minimum $N_{\max}$ and $L_{\max}$, since the number of coefficients is $N_{\max} \times L_{\max}^2$.
    We approach this problem in two parts -- first $L_{\max}$, then $N_{\max}$ -- since each term encodes different information -- the angular and radial resolution, respectively.

    The angular component of the SCF density function is $\propto K_{nl} Y_{lm}(\theta, \phi)$ \citep[eq 2.24]{HO92}, where $K_{nl}$ is a constant.

    For spherical patches, the corresponding resolution is [citation needed]

    \begin{equation}\label{eq:angular_resolution}
        A_{\Omega}(\ell) \approx \frac{4\pi}{(\ell+1)^2} \; \rm{[sr]}
    \end{equation}

    The radial component of the SCF is not distributed equally in $r$. Examining at \citet[eq 2.24]{HO92}, SCF density function is $\propto K_{nl} C_{n}^{\alpha}(\zeta)$, where $C$ are the Gegenbauer polynomials. These are functions of $\zeta \equiv \frac{{r}/{a} - 1}{{r}/{a} + 1}$, not $r$. The Gegenbauer polynomial $C_{n}^{\alpha}(\zeta)$ has $n$ zeros distributed approximately evenly over $\zeta \in (-1, 1)$. Stated in terms of the inter-$\zeta$ distance, $\Delta\zeta(n) \simeq \frac{2}{n}$.

    Starting from $\zeta$:
    \begin{flalign}
        \zeta &\equiv \frac{r/a - 1}{r/a + 1} \\
        d \zeta &= \frac{2}{(\frac{r}{a}+1)^2} \frac{dr}{a}
    \end{flalign}

    Taking $d\zeta \approx \Delta\zeta$ and $dr \approx \Delta r$

    \begin{equation}\label{eq:radial_resolution}
        \frac{\Delta{r}}{a}(n) \approx \frac{(\frac{r}{a}+1)^2}{n}
    \end{equation}

    Therefore, the smallest resolvable volume (per cubic scale length) is

    \begin{align}
        \frac{V(r, n, \ell)}{a^3}
        &= \frac{4\pi}{(\ell+1)^2} \left(\frac{r}{a}\right)^2 \frac{\Delta{r}}{a} \\
        &= \frac{4\pi}{(\ell+1)^2} \left(\frac{r}{a}\right)^2 \frac{(\frac{r}{a}+1)^2}{n}  \label{eq:resolution}
    \end{align}

    Note that the resolution is $r$-dependent. At fixed $n, \ell$ lower $r$ has higher resolution.
    Therefore, the important consideration for $N_{\max}$ and $L_{\max}$ will often be the required $\Delta r$ and and $A_{\Omega}$ at $r_{\max}$, the largest relevant $r$.

% section choosing_nmax_lmax (end)


% -----------------------------------------------------------------------------
%                                BIBLIGRAPHY
% -----------------------------------------------------------------------------

\newpage
\setlength{\bibsep}{6pt}
\bibliographystyle{abbrv}
\bibliography{main}


%------------------------------------------------------------------------------
\end{document}


%%%%%%%%%%%%%%%%%%%%%%%%%%%%%%%%%%%%%%%%%%%%%%%%%%%%%%%%%%%%%%%%%%%%%%%%%%%%%%#